%==============================================================================
\section{Conclusão e Trabalhos Futuros}
\label{sec:conclusion}

Colocamos uma camada de medição informacional sobre uma plataforma de
\emph{burn-in} de envelhecimento acelerado já existente, reexpressando a
contaminação ambiental, a correção PVT e a recuperabilidade do
envelhecimento como entropia, informação mútua e capacidade de canal. Em
uma campanha PID de \valDurationHours~h estabelecemos as âncoras empíricas
($\Ent{\slk}=\valHs$~bits; $R^2$ linear $=\valRsqLinear$) e uma metodologia
que substitui o $R^2$ dependente de modelo pela informação mútua, define
qualidade de \emph{burn-in} como a minimização de
$\MI{\slk}{\Tdie,\Vcore}$, converte o ruído medido em estados de
envelhecimento resolvíveis e uma capacidade de canal de alarme, e limita a
detectabilidade de RUL honestamente via Cram\'er--Rao e estimação
Bayesiana. Uma extensão via processo de Wiener e filtro de Kalman do mesmo
ferramental de estimação transforma ainda mais esse limite pontual de RUL
em uma trajetória de envelhecimento denoised, continuamente atualizada, e
uma \emph{distribuição} de RUL Gaussiana-inversa em forma fechada
(\cref{sec:res-wiener}), validada cruzadamente contra a estimativa
CRLB/Bayesiana e contra o resultado de seleção de modelo por MDL da
\cref{sec:res-correction}.

A proposta original deste projeto previa uma estimativa de capacidade
\emph{por regime térmico} (\emph{bang-bang} vs.\ PID); sem um registro bruto
de \emph{bang-bang} sincronizado e comparável ao da campanha PID de
\valDurationHours~h, essa comparação por regime não foi entregue como tal
--- em seu lugar, este trabalho entrega a capacidade e o limite de RUL no
regime PID (\cref{sec:res-capacity,sec:res-rul}) e, como sucedâneo direto
para a comparação entre regimes, a divergência distribucional KL/JS entre
as duas bancadas (\cref{sec:res-kl}), que já mostra numericamente que a
bancada não-PID é mais ruidosa. Completar a comparação de capacidade por
regime com um registro \emph{bang-bang} novo é, por isso, o primeiro item
de trabalho futuro abaixo, não uma lacuna silenciosa.

\paragraph{Trabalhos futuros.} Adquirir uma campanha \emph{bang-bang} de
duração equivalente para completar a comparação de KL e capacidade
entre regimes com dados novos; estender para múltiplos dispositivos para
separar uma trajetória de envelhecimento permanente da variação
dispositivo-a-dispositivo; implantar o \emph{pipeline} em janelas
deslizantes, on-line, como um gerador contínuo de atributos informacionais
alimentando o controlador de SLM, com a trajetória suavizada por Kalman da
\cref{sec:res-wiener} como sua estimativa de envelhecimento ao vivo;
calibrar um limiar de falha real $D$ (por exemplo, a partir de uma
especificação de margem de tempo) para que a \cref{eq:invgauss} produza uma
distribuição de RUL operacional, em vez dos limiares auto-referenciados
usados aqui; e prosseguir com a rota de ICA/separação cega com canais
derivados adicionais para fortalecer a recuperação de envelhecimento livre
de modelo.

\section*{Reprodutibilidade}
Todo escalar neste trabalho é gerado pelo \emph{pipeline} de análise em
\texttt{generated/values.tex}, e cada figura/tabela em \texttt{figures/} e
\texttt{tables\_pt/}; reexecutar a análise sobre um novo registro
regenera o documento sem transcrição manual.

\section*{Agradecimentos}
Este trabalho foi desenvolvido para a disciplina de Teoria da Informação
(TIP7266, PPGETI/UFC, 2026.1) e se apoia na pesquisa do autor sobre sensores
de envelhecimento de circuitos integrados e instrumentação de \emph{burn-in}.
